\chapter{Experimente și analiza rezultatelor}

\subsection{Analiza Arhitecturilor de Clasificare în CSNN}

Alegerea stratului final de decizie într-o rețea neuronală convoluțională de tip spiking (CSNN) reprezintă un compromis între eficiența computațională, precizia clasificării și plauzibilitatea biologică. În această secțiune, analizăm cele două abordări fundamentale utilizate în literatura de specialitate.

\subsubsection{Abordarea Hibridă: CSNN cu Clasificator SVM}

Această arhitectură utilizează straturile convoluționale de tip spiking exclusiv pentru extracția de trăsături (feature extraction), în timp ce decizia finală este delegată unui algoritm Support Vector Machine (SVM).

\textbf{Motivația utilizării:}
\begin{itemize}
    \item \textbf{Separabilitatea Datelor:} Trăsăturile extrase prin procese nesupervizate (precum STDP - Spike-Timing-Dependent Plasticity) nu sunt întotdeauna liniar separabile. SVM utilizează un kernel pentru a proiecta aceste date într-un spațiu de dimensiune superioară unde se poate determina un hiperplan optim de separare:
    \[ f(x) = \text{sgn} \left( \sum_{i=1}^{n} \alpha_i y_i K(x_i, x) + b \right) \]
    \item \textbf{Robustete pe Seturi de Date Reduse:} SVM implementează principiul minimizării riscului structural, fiind mult mai eficient decât un strat neuronal atunci când numărul de mostre de antrenament pentru anumite emoții este limitat.
    \item \textbf{Decuplarea Învățării:} Permite antrenarea rapidă a clasificatorului fără a fi necesară recalcularea întregii rețele spiking, facilitând experimentarea cu diferite meta-argumente de clasificare.
\end{itemize}

\subsubsection{Abordarea End-to-End: CSNN cu Straturi Fully Connected Spiking}

În această configurație, rețeaua rămâne integral în domeniul spiking, unde ultimele straturi sunt compuse din neuroni complet conectați (FC) care emit impulsuri pentru a reprezenta apartenența la o clasă.

\textbf{Motivația utilizării:}
\begin{itemize}
    \item \textbf{Eficiență Energetică și Hardware Neuromorfic:} O rețea pură poate fi implementată direct pe cipuri neuromorfice (ex. Intel Loihi, IBM NorthPole). Introducerea unui SVM ar necesita conversia semnalului înapoi în domeniul analogic/numeric, anulând avantajele de consum redus ale SNN.
    \item \textbf{Păstrarea Informației Temporale:} Spre deosebire de SVM, care procesează adesea rate de activare (spike rates) statice, straturile FC spiking pot utiliza codificarea temporală (Time-to-First-Spike). Dinamica membranei neuronale este descrisă de ecuația:
    \[ \tau_m \frac{dv}{dt} = -(v - v_{rest}) + R \cdot I(t) \]
    unde decizia este influențată direct de momentul precis al sosirii impulsurilor.
    \item \textbf{Antrenament prin Surrogate Gradients:} Utilizarea algoritmilor de backpropagation prin timp (BPTT) cu gradienți substituiți permite optimizarea întregului lanț de procesare, de la pixeli până la eticheta emoției, asigurând o coerență matematică a întregului sistem.
\end{itemize}

\section{Setul de date utilizat}
\section{Metodologia de testare a modelului}
\subsection{Cazul A: Recunoaștere pe imagini brute (cu fundal)}
\subsection{Cazul B: Recunoaștere pe imagini cu segmentare ROI}
\section{Metrici de performanță}
\section{Analiza impactului zgomotului vizual asupra preciziei modelului}